\documentclass{beamer}
\usetheme{Madrid}
\usecolortheme{default}

\title{3D Axialmaschine --- Hub Surface}
\subtitle{Domain-Partition Algorithm for Turbine Surfaces}
\author{Domain-Partition 3D Project}
\date{\today}
\institute{Block-Structured Meshing Pipeline}

\begin{document}

\begin{frame}
  \titlepage
\end{frame}

\begin{frame}{T1\_9 Turbine Hub --- Overview}
  \begin{columns}
    \begin{column}{0.55\textwidth}
      \begin{itemize}
        \item T1\_9 turbine hub: cylindrical surface at constant radius
        \item Input: \texttt{T1\_9\_hub\_raw.stl} (raw triangulated surface)
        \item Goal: extract quadrilateral block structure for CFD meshing
      \end{itemize}
    \end{column}
    \begin{column}{0.45\textwidth}
      \centering
      \includegraphics[width=0.9\textwidth]{T1_9_hub.png}
      \tiny\textit{Placeholder: 3D hub visualization}
    \end{column}
  \end{columns}

  \vspace{0.5cm}
  \centering
  \footnotesize
  Pipeline: STL $\rightarrow$ unwrap $\rightarrow$ cross-field $\rightarrow$ singularities $\rightarrow$ streamlines $\rightarrow$ blocks
\end{frame}

\note{
  Speaker notes:
  This is the first slide of an 11-slide presentation on the domain-partition algorithm.
  The T1\_9 test case is a turbine runner with a cylindrical hub surface.
  The full pipeline takes the raw STL surface, unwraps it to 2D, builds a seam-aware cross-field,
  detects singularities, integrates streamlines, and finally produces quadrilateral blocks.
  This slide sets the stage by introducing the geometry and the overall objective.
}

\end{document}
